\begin{figure}[htbp]
	\centering
	\subfloat[][\emph{$|\vartheta_2|=\tfrac{1}{3}|\vartheta_1|$,
		$|OC_2|=3|OC_1|$: al tendere di $|\vartheta|\to 0$
		il centro si allontana all'infinito lungo la normale a $\ub_O$}]{%
		\begin{tikzpicture}[>=Stealth, line width=0.8pt, scale=1.0]
			
			\coordinate (O)  at (0,0);
			\coordinate (uO) at ({1.5*cos(30)},{1.5*sin(30)});
			\coordinate (C1) at (-0.500, 0.866);
			\coordinate (C2) at (-1.500, 2.598);
			
			% Retta normale tratteggiata con freccia verso infinito
			\draw[dashed, gray!60]
			({-0.3*cos(120)},{-0.3*sin(120)}) --
			({3.0*cos(120)}, {3.0*sin(120)});
			\draw[->, dashed, gray!60]
			({3.0*cos(120)},{3.0*sin(120)}) --
			({3.6*cos(120)},{3.6*sin(120)})
			node[above left=1pt] {\small$\vartheta\to 0$};
			
			% Vettore u(O)
			\draw[->, thick] (O) -- (uO) node[right=2pt] {$\ub_O$};
			
			% Angolo retto
			\draw[thin]
			({0.18*cos(30)},{0.18*sin(30)}) --
			({0.18*cos(30)+0.18*cos(120)},{0.18*sin(30)+0.18*sin(120)}) --
			({0.18*cos(120)},{0.18*sin(120)});
			
			% Rette da C1 e C2 verso estremo di u(O)
			\draw[gray!50, thin] (C1) -- (uO);
			\draw[gray!50, thin] (C2) -- (uO);
			
			% Arco vartheta1 in C1
			\draw[-{Stealth[bend]}, thick]
			({-0.5+0.45*cos(300)},{0.866+0.45*sin(300)})
			arc[start angle=300, end angle=356, radius=0.45];
			
			\node at ({-0.4+0.62*cos(328)},{0.786+0.62*sin(328)})
			{\small$\vartheta_1$};
			
			% Arco vartheta2 in C2
			\draw[->, thick] ({-1.5+0.7*cos(300)},{2.598+0.7*sin(300)})
			arc[start angle=300, end angle=327, radius=0.7];
			\node at ({-1.4+0.88*cos(313)},{2.598+0.98*sin(313)})
			{\small$\vartheta_2$};
			
			% Quote a sinistra
			\draw[<->, thin]
			({-0.55*cos(30)},  {-0.55*sin(30)}) --
			({-0.55*cos(30)-0.5},  {-0.55*sin(30)+0.866})
			node[midway, left=3pt] {\small$\dfrac{|\ub_O|}{|\vartheta_1|}$};
			\draw[<->, thin]
			({-1.6*cos(30)},  {-1.6*sin(30)}) --
			({-1.6*cos(30)-1.5},  {-1.6*sin(30)+2.598})
			node[midway, left=3pt] {\small$\dfrac{|\ub_O|}{|\vartheta_2|}$};
			
			% Punti
			\fill (O)  circle (2pt)   node[below right=2pt] {$O$};
			\fill (C1) circle (2.5pt) node[left=3pt] {$C_1$};
			\fill (C2) circle (2.5pt) node[left=3pt] {$C_2$};
			
	\end{tikzpicture}}%
	\hspace{0.04\textwidth}%
	\subfloat[][\emph{le coordinate $(x^*,y^*)$ di un punto $P$
		rispetto al riferimento locale in $C$ entrano direttamente
		nella formula degli spostamenti rigidi}]{%
		\begin{tikzpicture}[>=Stealth, line width=0.8pt, scale=1.0]
			
			\coordinate (O) at (0,0);
			\coordinate (C) at (-0.500, 0.866);
			
			% Assi x, y
			\draw[->, thin] (-1.5,0) -- (2.2,0) node[right] {$x$};
			\draw[->, thin] (0,-0.4) -- (0,2.5) node[above] {$y$};
			\node[below left=0.1pt] at (O) {$O$};
			
			% Vettore u(O)
			\draw[->, thick] (O) -- ({1.5*cos(30)},{1.5*sin(30)})
			node[right=2pt] {$\ub_O$};
			
			% Componenti tratteggiate
			\draw[dashed, gray!60]
			({1.5*cos(30)},0) -- ({1.5*cos(30)},{1.5*sin(30)});
			\draw[dashed, gray!60]
			(0,{1.5*sin(30)}) -- ({1.5*cos(30)},{1.5*sin(30)});
			\node[below=0.1pt] at ({1.5*cos(30)/2}, 0) {\small$u_O$};
			\node[left=0.1pt] at (0, {1.5*sin(30)/2}) {\small$v_O$};
			
			% Centro C nel secondo quadrante
			\fill[white] (C) circle (3pt);
			\draw[thick] (C) circle (3pt);
			\node[above left=2pt] at (C) {$C$};
			
			% Assi x*, y* da C (prolungati)
			\draw[->, thin, gray!50] (-1.4, 0.866) -- (0.3, 0.866)
			node[right] {\small$x^*$};
			\draw[->, thin, gray!50] (-0.5, 0.1) -- (-0.5, 2.0)
			node[above] {\small$y^*$};
			
			% Linee tratteggiate proiezioni di C
			\draw[dashed, gray!60] (-0.5, 0) -- (C);
			\draw[dashed, gray!60] (0, 0.866) -- (C);
			
			% Quote
			\draw[<->, thin] (0,-0.6) -- (-0.5,-0.6)
			node[midway, below] {\small$x_C{=}{-}\dfrac{v_O}{\vartheta}$};
			\draw[<->, thin] (-1.3,0) -- (-1.3,0.866)
			node[midway, left=2pt] {\small$y_C{=}\dfrac{u_O}{\vartheta}$};
			
			% Tacche sugli assi
			\draw[thin] (-0.5, 0.05) -- (-0.5,-0.05);
			\draw[thin] (0.05, 0.866) -- (-0.05, 0.866);
			
	\end{tikzpicture}}
	\caption{Centro di rotazione nel caso piano.
		\textbf{\textit{(a)}}~Costruzione geometrica:
		$C$ si ottiene ruotando $\ub_O$ di $90^\circ$ nel verso
		di $\vartheta$ e riscalando di $1/|\vartheta|$.
		\textbf{\textit{(b)}}~Riferimento locale con origine in $C$
		e assi $x^*$, $y^*$ paralleli a $x$, $y$: le coordinate
		$(x^*, y^*)$ di un punto $P$ entrano direttamente
		nella~\eqref{eq:spost_da_CR}.}
	\label{fig:centro_rotazione}
\end{figure}
